% ID: SEL-D0011-BB-E04
% ESTADO: cotejada. Segunda pasada visual completa y borrador compilado revisados.
% Fuente: Selectividad/matematicas_aplicadas/EXTRA_SUPLENTE_Modelo-1-Enunciados (Mat-Ap-2021).pdf
% SHA-256: 4a6c7d34ff35c3e7678c52970b4d4558cee2c35224b46c66035cb0bc8d5a7ec8
% Página: 1; bloque B; ejercicio original 4.

\begin{enumerate}[label=\alph*)]
  \item \textbf{(2 puntos)} Sea $f$ una función de la que sabemos que la
  gráfica de su derivada, $f'$, es una parábola con vértice en el punto $(0,8)$
  que corta al eje de abscisas en los puntos $(-4,0)$ y $(4,0)$.
  \begin{enumerate}[label=\arabic*.]
    \item Dibuje la gráfica de $f'$.
    \item A partir de dicha gráfica, halle los intervalos de crecimiento
    y decrecimiento de $f$, así como las abscisas de los extremos relativos de $f$.
    \item Sabiendo que la gráfica de $f$ pasa por el origen de coordenadas,
    calcule la recta tangente a la gráfica de $f$ en el punto de abscisa $x=0$.
  \end{enumerate}
  \item \textbf{(0.5 puntos)} Calcule la derivada de la función
  $g(x)=(-3+x^2)\cdot e^{2x-1}$
\end{enumerate}
